% Iteraties en fractalen. Haeckel: een siphonofoor is een kolonie van steeds dezelfde bouwstenen, en haar tentakels
% vertakken zich op elke schaal opnieuw. Haeckel tekende in 1904 al de fractale vormen van de natuur.
% Motief: de bovenrand van de band is de Koch-kromme, van links naar rechts na 0, 1, 2 en 3 iteraties: telkens wordt
% elk lijnstuk vervangen door vier lijnstukken van een derde zo lang (vier stukken van 5,25 cm, basis op 10,5 cm,
% top op 12,02 cm). De punten zijn berekend met een kort Python-script.
\def\omslagbestand{omslag-siphonophoren.jpg}
\def\omslagbron{Ernst Haeckel, \textit{Siphonophorae}, uit \textit{Kunstformen der Natur}, 1904 (publiek domein)}
\def\omslagkoch{(0.000,10.500) -- (5.250,10.500) -- (7.000,10.500) -- (7.875,12.016) -- (8.750,10.500) -- (10.500,10.500) -- (11.083,10.500) -- (11.375,11.005) -- (11.667,10.500) -- (12.250,10.500) -- (12.542,11.005) -- (12.250,11.510) -- (12.833,11.510) -- (13.125,12.016) -- (13.417,11.510) -- (14.000,11.510) -- (13.708,11.005) -- (14.000,10.500) -- (14.583,10.500) -- (14.875,11.005) -- (15.167,10.500) -- (15.750,10.500) -- (15.944,10.500) -- (16.042,10.668) -- (16.139,10.500) -- (16.333,10.500) -- (16.431,10.668) -- (16.333,10.837) -- (16.528,10.837) -- (16.625,11.005) -- (16.722,10.837) -- (16.917,10.837) -- (16.819,10.668) -- (16.917,10.500) -- (17.111,10.500) -- (17.208,10.668) -- (17.306,10.500) -- (17.500,10.500) -- (17.597,10.668) -- (17.500,10.837) -- (17.694,10.837) -- (17.792,11.005) -- (17.694,11.174) -- (17.500,11.174) -- (17.597,11.342) -- (17.500,11.510) -- (17.694,11.510) -- (17.792,11.679) -- (17.889,11.510) -- (18.083,11.510) -- (18.181,11.679) -- (18.083,11.847) -- (18.278,11.847) -- (18.375,12.016) -- (18.472,11.847) -- (18.667,11.847) -- (18.569,11.679) -- (18.667,11.510) -- (18.861,11.510) -- (18.958,11.679) -- (19.056,11.510) -- (19.250,11.510) -- (19.153,11.342) -- (19.250,11.174) -- (19.056,11.174) -- (18.958,11.005) -- (19.056,10.837) -- (19.250,10.837) -- (19.153,10.668) -- (19.250,10.500) -- (19.444,10.500) -- (19.542,10.668) -- (19.639,10.500) -- (19.833,10.500) -- (19.931,10.668) -- (19.833,10.837) -- (20.028,10.837) -- (20.125,11.005) -- (20.222,10.837) -- (20.417,10.837) -- (20.319,10.668) -- (20.417,10.500) -- (20.611,10.500) -- (20.708,10.668) -- (20.806,10.500) -- (21.000,10.500)}
\def\omslagband{\fill[accent] (0,0) -- \omslagkoch -- (21,0) -- cycle;}
\def\omslagmotief{%
  \draw[white,line width=1.6pt,line join=round,line cap=round] \omslagkoch;
}
